\section{Basic objects}

\begin{definition}[Continuous alternating forms]\label{def:weak-form}
Let $X$ be a real normed space. A \emph{continuous alternating form} on $X$ is a continuous
bilinear map $b:X\times X\to\R$ such that $b(x,x)=0$ for every $x\in X$. Its induced operator is
\[
  S_b:X\longrightarrow\strongdual X,
  \qquad (S_bx)(y)=b(x,y).
\]
For a subspace $E\subseteq X$, define
\[
  E^{\circ_b}=\{x\in X: b(x,e)=0\text{ for every }e\in E\},
  \qquad
  \rad(b|_E)=E\cap E^{\circ_b}.
\]
\lean{KaltonPeck.Support.ContinuousAlternatingForm, KaltonPeck.Support.ContinuousAlternatingForm.orthogonal, KaltonPeck.Support.ContinuousAlternatingForm.radical, KaltonPeck.Support.ContinuousAlternatingForm.radical_eq_ker, KaltonPeck.Support.ContinuousAlternatingForm.radical_isClosed, KaltonPeck.Support.ContinuousAlternatingForm.restrict, KaltonPeck.Support.ContinuousAlternatingForm.restrictedRadical, KaltonPeck.Support.StrongSymplecticForm.toContinuousAlternatingForm}
\end{definition}

\begin{definition}[Strong symplectic form]\label{def:strong-form}
A continuous alternating form $\omega$ on a real normed space $X$ is \emph{strong} if its induced map
$D=S_\omega:X\to\strongdual X$ is a continuous linear isomorphism. A Banach space equipped with
such a form is called a strong symplectic Banach space.
\uses{def:weak-form}
\lean{KaltonPeck.StrongSymplecticForm, KaltonPeck.Support.StrongSymplecticForm}
\leanok
\end{definition}

\begin{definition}[Transpose and symplectic adjoint]\label{def:transpose-adjoint}
For a bounded operator $B:X\to Y$, its transpose is the bounded map
$B^*:\strongdual Y\to\strongdual X$ given by $B^*f=f\circ B$. If $X$ has a strong symplectic form with
induced isomorphism $D$ and $A:X\to X$ is a bounded endomorphism, define
\[
  A^+=D^{-1}A^*D.
\]
Equivalently, $\omega(A^+x,y)=\omega(x,Ay)$ for all $x,y\in X$.
\uses{def:strong-form}
\lean{KaltonPeck.transpose, KaltonPeck.Support.transpose, KaltonPeck.StrongSymplecticForm.adjoint, KaltonPeck.Support.StrongSymplecticForm.adjoint}
\leanok
\end{definition}

\begin{definition}[Fredholmness, rank, and nullity]\label{def:fredholm-rank}
A bounded map $A:X\to Y$ is \emph{Fredholm} when $\ker A$ is finite-dimensional,
$\ran A$ is closed, and $Y/\ran A$ is finite-dimensional. It has \emph{finite rank} when
$\ran A$ is finite-dimensional. In the respective finite-dimensional situations, set
\[
  \rank A=\dim\ran A,
  \qquad
  \operatorname{null}A=\dim\ker A,
  \qquad
  \operatorname{ind}A=\dim\ker A-\dim(Y/\ran A)\in\mathbb Z.
\]
\lean{KaltonPeck.IsFredholm, KaltonPeck.Support.IsFredholm, KaltonPeck.HasFiniteRank, KaltonPeck.Support.HasFiniteRank, KaltonPeck.operatorRank, KaltonPeck.Support.operatorRank, KaltonPeck.nullity, KaltonPeck.Support.nullity, KaltonPeck.Support.fredholmIndex}
\leanok
\end{definition}

\begin{definition}[Complex structures and hyperplanes]\label{def:complex-hyperplane}
A \emph{complex structure} on a real Banach space $X$ is a bounded operator $J:X\to X$ satisfying
$J^2=-I$. A \emph{hyperplane} is a closed linear subspace $H\subseteq X$ with
$\dim(X/H)=1$.
\lean{KaltonPeck.IsComplexStructure, KaltonPeck.Support.IsComplexStructure, KaltonPeck.IsHyperplane, KaltonPeck.Support.IsHyperplane}
\leanok
\end{definition}
